\documentclass{syllabus}

\begin{document}

\thispagestyle{fancy}
\date{Fall 2025}
\author{}
\title{PHY 3300: Thermal Physics \\ {\large \mdseries Classical Thermodynamics and Statistical Mechanics}}
\maketitle



\section*{General Information}

\raisebox{0.9\height}{

\begin{tabular}{lll}
\textbf{Instructor:}    & \multicolumn{2}{l}{Joseph Anderson}                             \\
\textbf{}               & \multicolumn{2}{l}{\href{mailto:janderson19@carthage.edu}{janderson19@carthage.edu}}                    \\
\textbf{}               & \multicolumn{2}{l}{DSC 082}                    \\
\textbf{}               & \multicolumn{2}{l}{(262)551-5871}                               \\ \\
\textbf{Office hours:}  & Wed, Fri               & 10:20 AM - 12:00 PM                    \\
\textbf{}               & Tuesday                & 3:30-5:00 PM (Discussion)              \\ \\
\textbf{Textbook:}      & \multicolumn{2}{l}{\textit{An Introduction to Thermal Physics}} \\
\textbf{}               & \multicolumn{2}{l}{Schroeder (Oxford 2021 or Longman 2000)}     \\ \\
\textbf{Prerequisites:} & \multicolumn{2}{l}{PHY 2210 (gen. phys. II)}                    \\
\textbf{}               & \multicolumn{2}{l}{Co-enrolled MTH 2020 (dif. eq.)}            
\end{tabular}
%
}
\hspace{\fill}
\includegraphics[height=2.25in]{Assets/thermalphysicsSchroederCover.jpg}
\hspace{.25in}

\section*{Course overview}

Thermal physics is a study of mechanics in the limit of large systems. New concepts come into play in these large systems which might sound familiar: temperature, entropy, transport, etc. Because the fundamental bodies of mechanics (e.g. atoms) are so small, nearly everything we might study can be considered a ‘large system,’ be it the motion of fluids in a single living cell, the gas dynamics in a combustion engine, or the evolution of galaxies and even the universe. Thermal physics will give us a common language and common principles by which we can reason about the behavior of all these systems and more.

Moreover, thermal physics can be studied in two ways: as its own science using principles deduced from experience, or by studying classical mechanics in the limit of many interacting bodies. The former is classical thermodynamics, the spark of the industrial revolution. The latter is statistical mechanics, a subject which not uncommonly drove men mad but also produced the transistor revolution. These two perspectives on thermal physics offer a rare combination of technological promise and existential conundrums. Both fundamentally ask similar questions: What does it mean that some processes can’t be reversed? How does energy distribute in a system? Why do objects behave so radically different at different temperatures? At the end of the course you will have the tools to tackle these questions and more.
\clearpage

\subsection*{Student learning outcomes} 

\noindent
Upon successful completion of this course, the student will be able to:
\begin{enumerate}[label=\textcolor{CarthageDark}{\bfseries\arabic*.}]
    \item Discuss the relationship between temperature, energy, and entropy, and apply these concepts to predict real-world behavior. 
    \item Analyze the distinction between reversible and irreversible processes and apply these distinctions to the development of heat engines.
    \item Apply fundamental probability theory to derive thermodynamic principles in classical and quantum mechanical systems.
    \item Compare and contrast the statistical mechanics of classical and quantum mechanical systems.
    \item Identify and justify the fundamental assumptions of thermodynamics and statistical mechanics.
    \item Analyze real-world datasets using principles of thermal physics.
\end{enumerate}

\section*{Course components}

\subsection*{Attendance}
This will be an active classroom: your participation is essential not only to your own learning but also to the learning of your peers. Your discussion and your interests will shape the direction of this course, but this can only occur if you first attend the class sessions. Attendance/participation will not be assessed as part of the course grade. However, this is subject to revision if attendance becomes an issue.

\subsection{Homework} 
Homework assignments will consist of a few problems from the text and will be due every Wednesday. Each week two students will have five minutes to present a homework problem. One fifth of the homework grade is based off of these presentations, for which a rubric will be supplied. These short explanations should focus on the key principle of the homework problem as well as enough explanation of the mathematical concepts to understand the solution. 

\subsection{Quizzes} Short quizzes will be administered once a week. These will consist of a brief, conceptual, written response problem. 

\subsection{Exams} There will be two in-class exams and a comprehensive final exam. Each exam is equally weighted. 


\subsection{Case studies}
There will be 3 case studies in this course drawn from current or historical events and important applications of the material in various related fields. These case studies should be worked on both in class and outside of class and require a written solution in the form of a short (3-5 page) paper with sufficient physical and mathematical insights to demonstrate your understanding of the material covered. 

\subsection{Reflection}
As part of an initiative in developing an \emph{entrepreneurial mindset} (EM), you will be keeping a reflection journal on EM behaviors. The purpose of this exercise is to help you consider how your actions in this class (or beyond) are contributing to your professional development.

\section*{Grading Scheme}
\subsection*{Component Weights}

\noindent\begin{tabular}{lr}
\textbf{Component} & \textbf{Weight}       \\ \hline
\textbf{Participation}            & 0\%  \\
\textbf{Homework}            & 12\%       \\
\textbf{Homework presentations}            &     3\%   \\
\textbf{Quizzes}            & 10\%       \\
\textbf{Exams} & 57\% \\
\textbf{Case studies}            & 15\%\\
\textbf{Reflection} & 3\%
\end{tabular}%



\lettergrades{a) number of extra credit assignments adequately completed, b) a pattern of seeking help in the course at office hours, c) number of missing homework assignments.}

\tipsforsuccess{upper}

\section*{Policies}
\subsection{General course policies}
Expect a response within one business day for emails sent to the instructor. To ensure a prompt response, format the subject line as: [PHY 3300] *insert subject here*

All required material will be covered in lecture, but can also be found in the parallel readings, so in the case of absence students should consult the scheduled readings. Students found sleeping will be brought to the attention of the class for public ridicule and no effort made to wake them up.

Students are expected to arrive and leave on time. I will attempt to arrive five minutes early and be available for at least five minutes after class.

\latework{Easy}
\collegepolicies

\end{document}
